\documentclass[11pt]{article}
\usepackage[margin=1in]{geometry}
\usepackage{booktabs}
\usepackage{enumitem}
\setlist{itemsep=2pt,topsep=3pt}

\title{CS 623 Assignment 2}
\author{Roll No.\ 230217}
\date{}

\begin{document}
\maketitle

\section*{Setup}
All numbers in this report were measured on the GPU3 department server:
NVIDIA A40 (compute capability 8.6, 46~GB), CUDA 12.6, built with
\texttt{nvcc -O2 -std=c++17 -arch=sm\_86} (Problem~2 adds \texttt{-fmad=false}
so that the GPU sums round exactly like the CPU code; Problem~3 adds
\texttt{-lineinfo} for the profiler). The GPU was shared: another user's job
kept it busy (about 90\% utilisation in \texttt{nvidia-smi}) during the runs,
so timings depend on what else is running. I therefore repeated each
measurement several times at different times of day and report the \emph{best}
time (contention can only add time) together with the worst time seen. For
the interleaved Problem~2 runs I logged \texttt{nvidia-smi} before each run, and
the best times were reproduced in runs that started with no other process on
the GPU. All CUDA allocations, copies, memsets, kernel
launches and synchronisations in Problems~2 and~3 are checked with the
\texttt{cudaCheckError} macro of the Problem~1 template (launches through
\texttt{cudaGetLastError} and the following synchronising call).

\paragraph{Build and run.} From the \texttt{230217-assign2} directory:
\begin{verbatim}
make                     # problem1, problem2 (v1,v2,v3), problem3 (2 variants)
./problem1/problem1.out
cd problem2 && ./problem2-v1.out && ./problem2-v2.out && ./problem2-v3.out
cd ../problem3 && ./problem3-vanilla.out N $PWD/insert_100000.bin \
                  $PWD/search_100000.bin $PWD/delete_100000.bin
\end{verbatim}
Problem~1 takes its size from the constant \texttt{Input\_Width} in
\texttt{230217-prob1.cu} (\texttt{1 << 6}, \texttt{1 << 7} or \texttt{1 << 8});
change it and rebuild with \texttt{make problem1}. Problem~2 must be run from
inside \texttt{problem2/}, because it opens \texttt{disp.txt} and
\texttt{grid.txt} from the current directory. The three versions of Problem~2
are one source file, \texttt{230217-prob2.cu}, built with
\texttt{-DVERSION=1}, \texttt{2} and \texttt{3}; the two variants of Problem~3
are one source file built with and without \texttt{-DBLOCKED}.

%------------------------------------------------------------------
\section*{Problem 1: 3D convolution}

\subsection*{Implementation}
The input is an $N\times N\times N$ array of \texttt{float}, the filter is a
$5\times5\times5$ cube, and cells outside the array are treated as zero.
Every version computes
\[
\mathrm{out}[z][y][x]=\sum_{f_z,f_y,f_x=0}^{4}
 \mathrm{filter}[f_z][f_y][f_x]\cdot \mathrm{in}[z+f_z-2][y+f_y-2][x+f_x-2].
\]
\begin{itemize}
\item \textbf{CPU:} single-threaded reference with six nested loops.
\item \textbf{Basic kernel:} one thread per output cell, $8\times8\times8$
 blocks and an $(N/8)^3$ grid. The input and the filter are read from global
 memory, up to 125 reads of each per thread.
\item \textbf{Optimised kernel:} the same mapping with
 (a) a $12\times12\times12$ shared-memory tile per block ($8+2\cdot 2$ in each
 dimension, ghost cells set to zero), loaded cooperatively by the 512 threads
 followed by \texttt{\_\_syncthreads()};
 (b) the filter in \texttt{\_\_constant\_\_} memory;
 (c) the innermost filter loop unrolled by hand.
\end{itemize}
Each kernel is run once as a warm-up and then 10 times; the program reports
the average. \emph{Kernel time} uses CUDA events around the launch;
\emph{end-to-end time} uses \texttt{chrono} around the host-to-device copy,
the kernel and the device-to-host copy. \texttt{check\_result} compares every
GPU output with the CPU output: it reported ``No differences found'' for both
kernels at all three sizes, both with the template's all-ones input and with
random integer input (so that indexing errors cannot hide behind a constant
input).

\subsection*{Results}
Times in ms. Each entry is the best of @@P1RUNS@@ program runs (made in two
sessions); every program run already averages 10 kernel launches. The ranges
over all runs are given in the second table.

\begin{center}
\begin{tabular}{rrrrrr}
\toprule
$N$ & CPU & Basic kernel & Basic end-to-end & Opt kernel & Opt end-to-end \\
\midrule
64  & @@CPU64@@  & @@BK64@@  & @@BE64@@  & @@OK64@@  & @@OE64@@ \\
128 & @@CPU128@@ & @@BK128@@ & @@BE128@@ & @@OK128@@ & @@OE128@@ \\
256 & @@CPU256@@ & @@BK256@@ & @@BE256@@ & @@OK256@@ & @@OE256@@ \\
\bottomrule
\end{tabular}
\end{center}

\begin{center}
\small
\begin{tabular}{rrrrrr}
\toprule
\multicolumn{6}{c}{Range over all runs (ms)} \\
$N$ & CPU & Basic kernel & Basic end-to-end & Opt kernel & Opt end-to-end \\
\midrule
64  & @@RCPU64@@  & @@RBK64@@  & @@RBE64@@  & @@ROK64@@  & @@ROE64@@ \\
128 & @@RCPU128@@ & @@RBK128@@ & @@RBE128@@ & @@ROK128@@ & @@ROE128@@ \\
256 & @@RCPU256@@ & @@RBK256@@ & @@RBE256@@ & @@ROK256@@ & @@ROE256@@ \\
\bottomrule
\end{tabular}
\end{center}

\begin{center}
\begin{tabular}{rrrr}
\toprule
$N$ & Basic kernel vs CPU & Opt kernel vs CPU & Opt kernel vs basic kernel \\
\midrule
64  & @@SPB64@@$\times$  & @@SPO64@@$\times$  & @@OVB64@@$\times$ \\
128 & @@SPB128@@$\times$ & @@SPO128@@$\times$ & @@OVB128@@$\times$ \\
256 & @@SPB256@@$\times$ & @@SPO256@@$\times$ & @@OVB256@@$\times$ \\
\bottomrule
\end{tabular}
\end{center}

\subsection*{Discussion}
In the basic kernel every input value is fetched from global memory by up to
125 different threads, and every thread also reads the same 125 filter weights
from global memory. The optimised kernel removes this redundancy: a block
computes an $8^3$ output chunk from a $12^3$ input region that is loaded from
global memory once ($1728$ reads instead of about $512\times125$) and then
reused from shared memory, the filter reads become constant-cache broadcasts
(all threads of a warp read the same weight at the same time), and the
unrolled inner loop gives the compiler more independent loads and multiplies
to schedule. The measured gain of the optimised kernel over the basic one
grows with the problem size, from @@OVB64@@$\times$ at $N=64$ to
@@OVB256@@$\times$ at $N=256$: at small sizes the kernel is so short
(tens of microseconds) that launch and scheduling overheads probably limit
both versions (I did not profile this), while at $N=256$ the memory traffic
saved by the tile matters.

The end-to-end numbers tell a different story. For $N=256$ the end-to-end
time of the basic kernel is about @@E2E256@@ times its kernel time, and the
optimised kernel's end-to-end time is only slightly lower than the basic one's
(the end-to-end times also vary the most between runs, see the table of ranges):
copying the $256^3$-float input and output over PCIe takes far longer than
either kernel, so a faster kernel barely changes the total unless the data
stays on the device across many launches.

%------------------------------------------------------------------
\section*{Problem 2: grid search}

\subsection*{Mapping the iteration space}
The grid has 10 variables with 13 values each ($0.5,\dots,6.5$ from
\texttt{grid.txt}; the count is computed as $\lfloor(\mathrm{end}-\mathrm{start})/\mathrm{step}\rfloor$ from
the file, never hard-coded), so $13^{10}\approx1.4\times10^{11}$ points. A
point is a result if $|\sum_j c_{ij}x_j-d_i|\le e_i$ for all 10 rows
($e_i=0.3\cdot ey_i$). A point is identified by one number in mixed radix,
\texttt{point\_id} $=((i_1 n_2+i_2)n_3+\dots)n_{10}+i_{10}$, stored as a 64-bit
integer because $13^{10}>2^{32}$. Ascending \texttt{point\_id} is exactly the
loop order of \texttt{problem2-v0.c}, so the matches are sorted on the host and
written in the same order and format as the reference. The coefficients,
bounds, steps and counts are copied to \texttt{\_\_constant\_\_} memory.

\subsection*{The three versions}
\begin{enumerate}
\item[\textbf{(i)}] \textbf{Simple port.} The host loops over variables 1 and
 2 (169 kernel launches); each launch has one thread per combination of
 variables 3--6 (28\,561 threads), and each thread loops over variables 7--10.
 All ten row sums are computed from scratch for every point. A thread that
 finds a match takes a slot with \texttt{atomicAdd} on a counter and stores
 the \texttt{point\_id}; the host reads the counter after each launch, copies
 the matches back and finally sorts them.
\item[\textbf{(ii)}] \textbf{Optimised.} The same exhaustive search with the
 basic optimisations of the cascaded partial-sum structure of my own earlier
 CPU version of this search: (a) \emph{reuse of partial sums} --- a table of running sums per row is
 kept, so changing the innermost variable costs one multiply-add per row
 instead of ten, and $d_i$ is folded into the first level; (b) \emph{early exit
 on the final check only} --- the last row test stops at the first row that
 fails; (c) \emph{fewer, larger kernel launches} --- the host loops only over
 variable 1 (13 launches) and each launch has $13^5=371\,293$ threads, which
 fills the GPU much better than the 28\,561 threads of (i). No point is
 skipped: every one of the $1.4\times10^{11}$ points is still checked.
\item[\textbf{(iii)}] \textbf{Thrust.} The point numbers are processed in
 slices of $2^{28}$. For each slice a \texttt{thrust::\allowbreak counting\_iterator} feeds
 \texttt{thrust::\allowbreak transform}, whose functor decodes the point number, evaluates
 the ten rows (returning 0 at the first failing row) and writes a flag;
 \texttt{thrust::\allowbreak count\_if} sizes the output and \texttt{thrust::\allowbreak copy\_if}
 with the flags as stencil collects the matching point numbers. The slices go
 in increasing order, so the output is already sorted.
\end{enumerate}
Shared-memory tiling and prefetching, which the assignment lists as examples,
are not used: the only data are the 120+30 numbers of the two input files
(about 1.2~KB), which already live in constant memory and are read with the
same address by all threads of a warp (a broadcast), so there is nothing to
tile; I did not measure them and make no claim about them.

\subsection*{Correctness}
@@V0LINE@@
All results use the same output format as \texttt{problem2-v0.c}. The files
were compared with \texttt{diff}:
\begin{center}
\begin{tabular}{lrl}
\toprule
Grid & Result points & Files compared \\
\midrule
6 values per variable ($6^{10}$ points)  & 2071  & v0 = v1 = v2 = v3 (byte-identical) \\
8 values per variable ($8^{10}$ points)  & 9316  & v0 = v1 = v2 = v3 (byte-identical) \\
13 values per variable (full grid)       & @@PNTS@@ & v0 = v1 = v2 = v3 (byte-identical) \\
\bottomrule
\end{tabular}
\end{center}

\subsection*{Performance}
Full grid, total time as printed by each program (host-side wall time of the
whole search: kernel launches or Thrust calls plus copying the matches back;
the final sort of the matches and the file output come afterwards and are not
included). Each version was run @@RUNS@@ times, interleaved with the other
versions, at different times while other users were also using the GPU.
\emph{Best} is the fastest run; it was reproduced in runs that started with
\texttt{nvidia-smi} showing no other process on the GPU (utilisation 0--14\%),
to within 2\% (e.g.\ version (ii) took 6.28~s in four separate runs). \emph{Worst} is the slowest run, when another
job was using the GPU at the same time. The CPU time is a single run of
\texttt{problem2-v0.c} on the server.

\begin{center}
\begin{tabular}{lrrrr}
\toprule
Version & Best (s) & Worst (s) & Speed-up vs (i) & Speed-up vs CPU \\
\midrule
CPU \texttt{problem2-v0.c} & @@T0@@ & --- & --- & 1 \\
(i) simple port    & @@T1@@ & @@W1@@ & 1.00 & @@VC1@@ \\
(ii) optimised     & @@T2@@ & @@W2@@ & @@SP2@@ & @@VC2@@ \\
(iii) Thrust       & @@T3@@ & @@W3@@ & @@SP3@@ & @@VC3@@ \\
\bottomrule
\end{tabular}
\end{center}

\subsection*{Discussion}
Version (ii) is clearly the fastest, @@SP2@@ times faster than (i): reusing the
partial sums cuts the work per point from about 100 multiply-adds to a handful,
the early exit avoids evaluating the remaining rows of a failing point (results
are rare: @@PNTS@@ of $1.4\times10^{11}$ points), and the larger launches keep
the GPU better occupied. Versions (i) and (iii) are about equal on an idle GPU
(@@T1@@~s and @@T3@@~s). In the Thrust version every point is an independent
element, so each element decodes its own 64-bit point number (ten 64-bit
divisions and remainders, which are expensive on a GPU) and no work can be
shared between neighbouring points, whereas (i) decodes a prefix once per
thread and then runs through $13^4$ points; on the other hand (iii) stops at
the first failing row and (i) does not. I did not profile either version, so
this is my reading of the code, not a measurement of why they come out equal.
Thrust gives short code, and \texttt{transform} plus a predicate is all the
problem needs, but its element-wise model does not expose the loop structure
that (ii) exploits. The effect of sharing the GPU is large: under load the
times rose to @@W1@@~s, @@W2@@~s and @@W3@@~s, and (i) and (iii) changed places
between runs, which is why only the idle-GPU best times are compared.

%------------------------------------------------------------------
\section*{Problem 3: counting Bloom filter}

\subsection*{Design}
\begin{itemize}
\item \textbf{Structure.} The filter is a raw array of $N$ 8-bit counters on the
 GPU, stored as $\lceil N/4\rceil$ 32-bit words. Counter \texttt{pos} is byte
 \texttt{pos \% 4} of word \texttt{pos / 4}; \texttt{insert}, \texttt{delete}
 and \texttt{search} are three separate kernels with one thread per key, run
 one after the other from the host, so they never overlap. $k=3$.
\item \textbf{Atomics on bytes.} CUDA has no 8-bit atomics, so an increment is
 \texttt{atomicAdd(word, 1 << (8*(pos\%4)))} (counter overflow is ignored, as
 allowed). A decrement uses an \texttt{atomicCAS} loop on the containing word
 that refuses to go below zero; a plain subtraction would borrow from the
 neighbouring counter when the byte is 0.
\item \textbf{Delete.} Delete runs in two phases within one batch. A search
 kernel first records, for every key of the batch, whether all three of its
 counters are above 0, using the counters as they are before any decrement;
 then a delete kernel decrements the three counters (with the \texttt{atomicCAS}
 loop, never below 0) of the keys that were flagged. A key whose counters show
 it is absent therefore changes nothing. Taking the decision for all keys
 before any decrement makes the result independent of the order in which
 threads run. A key that was never inserted but whose three counters are all
 set by other keys (a false positive) cannot be told apart from a member by
 the filter and is decremented like one; this is an inherent limit of counting
 Bloom filters (see the experiment below), and the main traces delete only
 keys that were inserted.
\item \textbf{Duplicates.} Each trace is sorted and de-duplicated on the host
 before it is copied to the GPU, so a key repeated within a trace increments
 its counters once. If a key were inserted in two different insert traces its
 counters would count both inserts; the test traces use one insert trace.
\item \textbf{Hash functions.} \texttt{hash1}\dots\texttt{hash4} are separate
 \texttt{\_\_host\_\_ \_\_device\_\_} functions so that the host test computes
 exactly the same positions as the kernels. Each is
 $\mathrm{scramble}(\mathit{key}\cdot m+s)$ with its own multiplier $m$ and
 seed $s$, where \texttt{scramble} xors the value with two rotations of
 itself (different rotation amounts per hash), multiplies by a second odd
 constant and xors with a right shift by 15. I checked them statistically
 (uniform over 128 and 1024 buckets, about 16 of 32 output bits flip per input
 bit flip, near-independent pairs).
\item \textbf{Vanilla variant.} The three counters are
 \texttt{hash1,2,3(key) \% N}.
\item \textbf{Blocked variant.} The array is cut into 128-byte blocks (one L1
 cache line); \texttt{hash4(key)} selects the block and \texttt{hash1,2,3(key)
 \% 128} select the three counters inside it, so the $k+1=4$ hash functions
 of the assignment are all used and all counters of a key lie in one cache
 line. $N$ is rounded up to a multiple of 128. (The assignment names the first
 hash as the block selector; I use \texttt{hash4}, which gives the same
 structure.) If two of a key's counters coincide, that counter is incremented
 (and later decremented) once per hash, so insert and delete stay consistent.
\end{itemize}

\subsection*{Traces and test routine}
\texttt{trace\_generator.cpp} writes three binary files of \texttt{uint32} keys
for a given $n$ (no random-number generator: key $i$ is a fixed invertible
scrambling of the number $i$, so keys for different numbers never coincide).
The insert trace has keys $0..n-1$ with every tenth entry a repeat; the delete
trace has every third inserted key; the search trace contains, in a shuffled
order, a third keys that stay in the filter, a third keys that are deleted, and
a third keys that were never inserted. The submission contains the traces for
$n=100\,000$ (110\,000 / 100\,000 / 33\,334 keys, under 0.5~MB each).

\texttt{check\_result()} keeps a host model: a plain byte array updated with the
same rules (and the same \texttt{positions()} function) plus a
\texttt{std::set} of the keys that should be present. After the insert and
after the delete it copies the GPU counters back and counts how many differ
from the model; after each search batch it compares every answer with the set
and counts false negatives and false positives. The false-positive rate is the
number of false positives divided by the number of all searches, as defined in
the assignment. The set is used only in this test, never by the kernels.

\subsection*{Results}
On these traces, every run: 0 counters differ from the host model after insert
and after delete, and 0 false negatives. Search trace of $n$ keys, a third of which were never
inserted (so the rate among the keys that really are absent is three times the
figure shown after the insert, and one and a half times after the delete,
when two thirds of the keys are absent).

\begin{center}
\begin{tabular}{lllrr}
\toprule
Trace & $N$ & Variant & FP after insert & FP after delete \\
\midrule
$n=10^5$ (100k inserted)  & $10^6$     & vanilla & 0.547\% & 0.389\% \\
                          &            & blocked & 0.662\% & 0.554\% \\
$n=2\cdot10^6$            & $2\cdot10^7$ & vanilla & 0.583\% & 0.390\% \\
                          &            & blocked & 0.665\% & 0.500\% \\
$n=10^5$                  & $3\cdot10^5$ & vanilla & 8.42\%  & 7.73\% \\
                          &            & blocked & 8.52\%  & 7.91\% \\
\bottomrule
\end{tabular}
\end{center}
The blocked variant has the higher false-positive rate in every row, as the
assignment predicts: keys that fall into the same 128-counter block compete
for those counters, and unevenly loaded blocks cost more false positives than
lightly loaded blocks save. The gap is largest when the filter is lightly
filled and shrinks as it fills up.

\subsection*{Deleting keys that were never inserted}
To test the limit described under \emph{Delete}, I generated a delete trace in
which 10\,020 of its 43\,354 keys were never inserted (generator option
\texttt{ghost\_pct}\,=\,30) and ran it with $N=3\cdot10^5$, so that many
such keys look present. The insert and search traces are the $n=10^5$ ones.
Each program was run three times; the three runs gave identical output.

\begin{center}
\begin{tabular}{lrrrr}
\toprule
Variant & Never-inserted keys & Counters differing & False negatives & FP rate \\
 & that looked present & from host model & after delete & after delete \\
\midrule
vanilla & 2469 & 0 (all 3 runs) & 2025 & 6.77\% \\
blocked & 2549 & 0 (all 3 runs) & 1989 & 6.90\% \\
\bottomrule
\end{tabular}
\end{center}
The GPU counters match the sequential host model exactly. The false negatives
are real members whose counters were decremented by never-inserted keys that
looked present; no counting Bloom filter can avoid this, because the filter
cannot tell such a key from a member. In an earlier version of the delete
kernel, where each thread checked its counters and decremented them without
waiting for the other threads, the same trace gave about 2\,650 counters that
differed from the host model, and a different number in every run (2657, 2659
and 2655); this is why delete now takes its decision for all keys of a batch
before any decrement.

\subsection*{Profiling with Nsight Compute}
Profiled with Nsight Compute 2026.1 on the A40, for the first launch of each
kernel on the $n=2\cdot10^6$ traces with $N=2\cdot10^7$ counters (2.2 million
keys inserted, 2 million searched). Both variants were profiled.

\begin{center}
\small
\begin{tabular}{lrrrr}
\toprule
 & \multicolumn{2}{c}{Vanilla} & \multicolumn{2}{c}{Blocked} \\
Metric & insert & search & insert & search \\
\midrule
@@PROFROWS@@
\bottomrule
\end{tabular}
\end{center}
The kernels execute no floating-point instructions, so the arithmetic
intensity in FLOP/byte is exactly zero; the integer-instruction rate per byte
is given as the nearest meaningful figure. Warp efficiency is the
\texttt{smsp\_\_thread\_inst\_executed\_per\_inst\_executed} ratio (32 means
every lane is active), which is expected because the kernels have no
data-dependent branches inside the warp except the final bounds check.

\paragraph{Top three bottlenecks reported by the profiler} (in order of the
profiler's estimated speed-up; it names the same causes for both variants).
\begin{enumerate}
\item \textbf{Uncoalesced global memory accesses.} Insert: 84\% of the
 sectors are excess (estimated speed-up 83\%, vanilla and blocked alike).
 Search: 83\% of the sectors are excess (estimated speed-up 82\% vanilla,
 81\% blocked), and the DRAM access-pattern check adds that only 5.1 of the 32
 bytes of each fetched sector are used by a thread, with 90.5\% (vanilla) /
 81.2\% (blocked) of those sectors missing in L2 (estimated speed-up 61\% /
 45\%). Every key hashes to unrelated places, so neighbouring threads touch
 different cache lines.
\item \textbf{Memory-bound execution with long-scoreboard stalls.} Memory is
 busier than compute: DRAM throughput is 68--85\% of peak, and a warp spends
 288 (vanilla) / 202 (blocked) cycles per search waiting on an L1/global load
 (estimated speed-up 20\% / 15\%). The insert kernel moves about twice the
 DRAM traffic of the search kernel, presumably because each atomic is a
 read-modify-write.
\item \textbf{Low issue rate / occupancy.} Each scheduler issues an
 instruction only every 35--146 cycles with 0.01--0.04 eligible warps per cycle,
 because warps wait on memory, and achieved occupancy (70--84\%) is below the
 theoretical 100\% (estimated speed-up 15--32\%).
\end{enumerate}
Atomics are not named as a bottleneck by the profiler for these runs; their
cost shows up in the insert kernel's higher DRAM traffic and longer duration.

\paragraph{Effect of blocking.} The blocked variant is about @@BLKINS@@$\times$ faster
for insert and @@BLKSRC@@$\times$ for search, moves 36--37\% less DRAM data, and has a
higher L2 hit rate (36\% / 20\% against about 10\%): the three counters of one
key share a cache line, so one 128-byte line serves all three accesses. The
price is the higher false-positive rate in the table above. Both kernels remain
limited by uncoalesced access, since different keys still go to unrelated
blocks.

\end{document}
